\begin{proposition}[P4.18]\hypertarget{tag-02WP}{}\label{night:P4.18}\label{stmt:P4.18}
\StatementRecord{night:P4.18}{proved}{02WP}{\SeriesRef{night:P4.15}, \SeriesRef{night:P4.16}}

For $a\geq2$ and $b,c\geq1$, the reduced germ
$x^a-y^bz^c=0$ cannot occur with a common positive multiplicity
in a torus degeneration with smooth total threefold. Its affine
nearby fibre contains a spherical class of nonzero square.
In particular this excludes $x^2-z^3y^4$, whose smooth
normalization has a divisorial differential rank drop.
\end{proposition}
\begin{proof}

\begin{enumerate}[label={$\langle1\rangle$\arabic*.},leftmargin=*]
\item Write $d=\gcd(b,c)$, $b=dB$, $c=dC$, so $\gcd(B,C)=1$.
The nearby fibre $F=\{x^a-y^bz^c=1\}$ maps to the smooth
affine curve $A=\{x^a-w^d=1\}$ by $w=y^Bz^C$.
The function $w$ on $A$ has $a$ simple zeros. Choose an
embedded arc between two zeros avoiding the others, and a
small open disc $D$ around it. In a holomorphic coordinate
$\xi$ on $D$, write $w=u(\xi)(\xi-\alpha)(\xi-\beta)$
with $u$ a unit. Taking a $B$-th root of $u$, rescaling $y$,
and making an affine coordinate change identifies the preimage
of $D$ with
\[
 W=\{y^Bz^C=X^2-1,\ X\in D'\},
\]
where $D'$ is a disc containing both $-1$ and $1$.

\item Choose $Bp+Cq=1$. With $h=X^2-1$, the explicit trivialization
away from the two special fibres is
\[
 y=h^pt^C,\quad z=h^qt^{-B},\quad t=y^qz^{-p}.
\]
Its fundamental group is $F(\beta_-,\beta_+)\times\mathbb Z\tau$.
The neighbourhood of each special fibre is the graph
$X=\pm\sqrt{1+y^Bz^C}$ on $|y^Bz^C|<\delta$, a contractible
domain by real scaling. Its punctured part has two generating
loops in $y,z$, whose images in $\langle\beta_\pm,\tau\rangle$
are the columns of
\[
 \begin{pmatrix}B&C\\q&-p\end{pmatrix},\qquad \det=-1.
\]
Thus adjoining each neighbourhood kills $\beta_\pm$ and $\tau$.
Van Kampen gives $\pi_1(W)=0$.

\item The map
\[
 \pi:W\longrightarrow W'=\{y^BZ=X^2-1,\ X\in D'\},
 \qquad (X,y,z)\longmapsto(X,y,z^C)
\]
is proper finite of degree $C$. The compact class $v$ in
P4.16, of square $-2B$, is already supported over $D'$.
Indeed, the cover just given with $C=1$ has the identical
Mayer--Vietoris maps before and after enlarging $D'$ to the
whole plane: both punctured bases have the same two meridians,
and the same two graph neighbourhoods are attached. The
inclusion is therefore an isomorphism on $H_2$ and preserves
the primitive class and its square.
For its compact-support dual $\alpha\in H_c^2(W';\mathbb Z)$,
properness and degree give
\[
 \int_W(\pi^*\alpha)^2=C\int_{W'}\alpha^2=-2BC.
\]
One may verify this by smooth compactly supported forms:
outside the measure-zero branch divisor the map is an oriented
$C$-sheeted covering. The dual of $\pi^*\alpha$ is spherical
by $\pi_1(W)=0$ and Hurewicz. Its compact representatives
embed in $F$ with the same nonzero square.

\item The polynomial has positive weights $(b+c,a,a)$. Weighted
scaling places these representatives in arbitrarily small
nearby fibres. P4.16's explicit unit removal handles the unit
and common multiplicity. The spherical class becomes zero
in a torus, contradicting its square $-2BC$.

\item For the mixed example the normalization is the finite
birational map
\[
 (u,v)\longmapsto(u^3v^2,v,u^2).
\]
Finiteness follows from $u^2=z$, and $u=x/(zy^2)$ gives
birationality. The image ring contains all even powers of $u$
and the odd monomials $u^iv^j$ with $i\geq3,j\geq2$.
Its conductor is $u^2v^2\mathbb C\{u,v\}$, giving divisor
$2\{u=0\}+2\{v=0\}$. Its differential minors are
$3u^2v^2,-4u^4v,-2u$, with common divisorial factor $u$.
Thus this example explicitly lies outside P4.15's finite
rank-drop assumption. The 64 exponent matrices for
$1\leq b,c\leq8$, transfer squares and normalization
identities are executed in
\texttt{code/part4/binomial\_milnor.py}; its output and
JSON data are saved alongside it.

\end{enumerate}
\end{proof}